\sscp{**kusay `bear cuscus'}{14-03-01}
As mentioned above,
two species of cuscus are found widely on Sulawesi,
the ‘bear cuscus, \it{Ailurops ursinus}’,
and the ‘Sulawesi dwarf cuscus, \it{Strigocuscus celebensis}’.
It is worth noting that the bear cuscus is reported to have musk
glands that were of value to Chinese merchants.
The bear cuscus thus became a trade commodity of some value.

Terms for ‘cuscus’ are available for 75 languages of Sulawesi.\footnote{
		We are grateful to David Mead for compiling the data on Sulawesi
		forms for cuscus and making it available to us.}
Of these, 56 are known to have a term cognate with **kusay.
Whenever more detailed information is available,
this term designates the ‘bear cuscus’,
and another term designates the ‘dwarf cuscus’.
The most common term for ‘dwarf cuscus’ is from **taŋali,
known for 23 languages in Sulawesi across different subgroups.
The reconstructions for reflexes of **kusay ‘bear cuscus’ across
Sulawesi are given in \trf{14-01}.

\begin{table}[h]
	\centering\tcp{Reconstructions relating to **kusay bear cuscus in Sulawesi}{14-01}
	\st{5.5pt}\begin{tabular}{l|llll}\lsptoprule
Proto language		&	Form		&	gloss				&	Node								&	Source	\\	\midrule
P.-Sangir					&	**kusay	&	cuscus			&	\tsc{Sangir}				&	\citet[89]{Sneddon1986}	\\
P.-Gorontalo-			&	**kusoy	&	bear cuscus	&	\tsc{Greater C.}		&	\citet[322]{Usup1986}	\\
\hp{P.-}Mongondow	&					&							&	 \tsc{Philippines}	&	\hp{\,}	\\
P.-Minahasan			&	**kuse	&	cuscus			&	\tsc{Minahasan}			&	\citet[151]{Sneddon1978}	\\
P.-Celebic				&	**kuse	&	bear cuscus	&	\tsc{Celebic}				&	Michael Martens,	\\
									&					&							&											&	David Mead	\\
P.-S. Sulawesi		&	**kuse	&	cuscus			&	\tsc{S. Sulawesi}		&	\citet[742]{Mills1975}	\\
\lspbottomrule
	\end{tabular}
\end{table}

Languages of Sulawesi with a cognate of **kusay belong to at
least three contiguous subgroups of Malayo-Polynesian: \tsc{Philippine}
(including \tsc{Sangir}, \tsc{Minahasan}, and \tsc{Mongondow-Gorontalo}), \tsc{South Sulawesi}, and \tsc{Celebic}.
While the \tsc{Philippine} node is controversial,\footnote{
		See \citet{Blust2019} for recent arguments in favour of \tsc{Philippine}
		as a distinct subgroup,
		with the debate continued by \citet{Ross2020,Reid2020,Zorc2020,Liao2020,Blust2020}.}
the other nodes are generally accepted.
Given that marsupials are not found west or north of Sulawesi
(not found north of Sangir-Talaud),
**kusay ‘(bear) cuscus’ cannot credibly be assigned to PMP.
Thus, unless one is willing to propose a supergroup which contains
all these subgroups, borrowing must have played at least some role in the distribution
of **kusay across Sulawesi.

The Sangir languages and two languages of \tsc{Mongondow-Gorontalo}
(\tsc{Greater Central Philippine}) in the north preserve a final vowel-glide
sequence within this form: Ponosakan \it{kusoy},
Mongondow \it{kutoy} (regular *s > \it{t}),
Bantik, Ratahan \it{kusey}, and Sangir \it{kusai, kuse}.
Note that two of these languages have /o/ in the final syllable,
similar to forms in the Papuan \tsc{North Halmahera} languages across the strait.
All other cognates of **kusay have final *ay > \it{e},
which is a regular sound change in \tsc{Celebic}
and \tsc{South Sulawesi} subgroups.
Banggai \it{kuai} ‘bear cuscus’ also preserves a final vowel-glide
sequence, but that, along with the loss of /s/ flag the Banggai form as irregular.